%*******************************************************
% 4.4 Đối tượng và phạm vi; 4.5 Bố cục luận án
%*******************************************************
\section{Đối tượng và phạm vi nghiên cứu}\label{sec:phamvi}

\begin{tomtat}
Đối tượng là khác biệt chất lượng thị trường giữa các sàn và tác động của việc đổi sàn. Phạm vi gồm ba mẫu ứng với ba câu hỏi, dữ liệu từ 6/2009 đến hết 24 tháng sau sự kiện. Đề tài nói rõ những gì nằm ngoài phạm vi: câu hỏi số 2 và số 4 của bài toán số 19, thị trường trái phiếu và phái sinh (trừ phần định hướng), các chính sách khác và khảo sát nhà đầu tư.
\end{tomtat}

Đối tượng nghiên cứu là khác biệt chất lượng thị trường giữa các sàn và tác động của việc đổi sàn lên thanh khoản và khả năng đầu tư, quan sát qua đợt hợp nhất cổ phiếu niêm yết của HNX vào HOSE, đặt trong khung năng lực cạnh tranh của trung tâm tài chính.

Phạm vi nghiên cứu (dự kiến):
\begin{itemize}
\item[+] Phạm vi không gian: các trung tâm tài chính có điểm GFCI là mẫu của câu hỏi 1; sàn niêm yết HNX, UPCoM và HOSE là mẫu của câu hỏi 2; các cổ phiếu chuyển từ HNX sang HOSE ngày 28/12/2026 và các nhóm đối chứng là mẫu của câu hỏi 3.
\item[+] Phạm vi thời gian: dữ liệu thị trường từ tháng 6/2009 đến hết 24 tháng sau sự kiện; dữ liệu GFCI từ kỳ 1 đến kỳ mới nhất tại thời điểm phân tích.
\item[+] Phạm vi nội dung: danh mục chỉ mua là cấu hình cơ sở của bảy mô hình phân bổ. Thị trường trái phiếu và phái sinh chỉ được xét trong phần định hướng chuyên môn hóa. Đề tài không đánh giá chính sách nào khác ngoài đợt chuyển sàn và không khảo sát nhà đầu tư.
\item[+] Giới hạn khác: kết luận về \HH{1} là liên hệ có điều kiện; kết luận về \HH{3} là tác động gộp của sàn trong một sự kiện, khi các giả định ở mục~\ref{sec:did} được thỏa mãn; kết luận về \HH{4} chỉ áp dụng cho rổ CS.
\end{itemize}

Ba giới hạn ở gạch đầu dòng cuối nên được nhắc lại mỗi khi trình bày kết quả. \HH{1} không nói thị trường cổ phiếu \emph{gây ra} vị thế trung tâm. \HH{3} không tách được biên độ, bước giá và chỉ số ra khỏi nhau, và chỉ đúng cho đợt chuyển sàn này. \HH{4} chỉ đo trên 30 mã thanh khoản nhất, nên không nói được gì về các mã nhỏ ít giao dịch.

Đề tài không giải quyết trọn bài toán số 19: câu hỏi số 2 và số 4 (về công nghệ tài chính và dịch vụ tài chính thế hệ mới) nằm ngoài phạm vi, và với câu hỏi số 3, đề tài chỉ cung cấp thước đo nền. Người dùng dự kiến của kết quả gồm Sở Khoa học và Công nghệ Hà Nội, cơ quan đầu mối theo dõi và đánh giá kết quả thực hiện Danh mục các bài toán lớn (Ủy ban nhân dân thành phố Hà Nội, 2026, Điều 3), Ủy ban nhân dân thành phố Hà Nội, Ủy ban Chứng khoán Nhà nước, HNX, HOSE và các tổ chức đầu tư.

Trước khi nộp hồ sơ đề nghị hỗ trợ, thí sinh sẽ xin văn bản xác nhận hoặc cam kết sử dụng kết quả của ít nhất một đơn vị trong số này, ưu tiên HNX; văn bản này cũng là căn cứ của chỉ tiêu (6).

HNX được ưu tiên vì là đơn vị chịu ảnh hưởng trực tiếp nhất của đợt hợp nhất và là đơn vị sẽ vận hành các thị trường mà phần định hướng chuyên môn hóa hướng tới. Chỉ số chất lượng thị trường vốn, nếu được mở rộng sang trái phiếu và phái sinh như nêu ở mục~9, có thể trở thành công cụ theo dõi của chính HNX.

%=======================================================
\section{Dự kiến bố cục các chương và các tiểu mục đến 3 chữ số}\label{sec:bocuc}
%=======================================================

Ngoài phần kết luận, danh mục tài liệu tham khảo và phụ lục, luận án dự kiến gồm sáu chương (Bảng~\ref{tab:bocuc}).

\begin{table}[H]
\myfloatalign
\caption{Dự kiến bố cục các chương và tiểu mục của luận án}\label{tab:bocuc}
{\bangnho
\begin{tabularx}{\textwidth}{@{}l >{\raggedright\arraybackslash}p{0.25\textwidth} L@{}}
\toprule
Chương & Tên chương & Các tiểu mục chính (đến ba chữ số)\\
\midrule
Chương 1 & Giới thiệu nghiên cứu & 1.1 Bối cảnh: Hà Nội, bài toán số 19 và đợt hợp nhất sàn năm 2026; 1.1.1 Trung tâm tài chính quốc tế của Việt Nam; 1.1.2 Lộ trình chuyển cổ phiếu niêm yết từ HNX sang HOSE; 1.2 Vấn đề nghiên cứu; 1.3 Mục tiêu và câu hỏi nghiên cứu; 1.4 Đối tượng và phạm vi; 1.5 Đóng góp dự kiến; 1.6 Kết cấu luận án\\
\addlinespace
Chương 2 & Tổng quan tài liệu và khung phân tích & 2.1 Tích tụ và phân vai của trung tâm tài chính; 2.1.1 Phân vai các trung tâm trong một quốc gia; 2.2 Chất lượng thị trường vốn; 2.3 Đổi sàn và hợp nhất sàn; 2.3.1 Bằng chứng từ Euronext; 2.4 Đa dạng hóa danh mục khi tham số được ước lượng; 2.5 Nghiên cứu về Việt Nam; 2.6 Khoảng trống, khung phân tích và giả thuyết\\
\addlinespace
Chương 3 & Dữ liệu và phương pháp & 3.1 Dữ liệu quốc tế; 3.2 Dữ liệu HNX, UPCoM và HOSE; 3.2.1 Cổ phiếu chuyển sàn và thiên lệch sống sót; 3.2.2 Giai đoạn căng thẳng; 3.3 Chỉ số chất lượng thị trường vốn; 3.4 Chỉ số khả năng đầu tư; 3.5 Thiết kế khác biệt kép quanh ngày 28/12/2026; 3.5.1 Nhóm đối chứng và kiểm tra xu hướng trước sự kiện; 3.6 Đăng ký trước và các lệch khỏi kế hoạch\\
\addlinespace
Chương 4 & Phát triển thị trường cổ phiếu và năng lực cạnh tranh của trung tâm tài chính: bằng chứng quốc tế & 4.1 Dữ liệu bảng các trung tâm tài chính; 4.1.1 Số hóa và chuẩn hóa điểm GFCI; 4.2 Kiểm định \HH{1}; 4.3 Kiểm tra độ vững; 4.4 Thảo luận: Thành phố Hồ Chí Minh, Đà Nẵng và Hà Nội\\
\addlinespace
Chương 5 & Chất lượng thị trường vốn trước hợp nhất: HNX, UPCoM và HOSE & 5.1 Các chỉ báo thành phần; 5.1.1 Ảnh hưởng của biên độ dao động giá; 5.2 Chỉ số theo sàn và kiểm định \HH{2}; 5.3 Giai đoạn căng thẳng (mô tả); 5.4 Thảo luận\\
\addlinespace
Chương 6 & Tác động của việc chuyển sàn và định hướng cho thị trường vốn Hà Nội & 6.1 Thanh khoản quanh ngày 28/12/2026 và kiểm định \HH{3}; 6.1.1 Khác biệt của tác động theo mức chịu cú sốc biên độ, bước giá và chỉ số; 6.1.2 Khác biệt theo quy mô và theo trụ sở doanh nghiệp; 6.2 Khả năng đầu tư trước và sau sự kiện và kiểm định \HH{4}; 6.3 Định hướng chuyên môn hóa cho thị trường vốn Hà Nội; 6.4 Đề xuất chính sách và hạ tầng dữ liệu\\
\addlinespace
Kết luận & & Trả lời ba câu hỏi nghiên cứu; hạn chế (một sự kiện duy nhất, độ dài dữ liệu sau sự kiện, các cú sốc trùng thời gian, \HH{1} là liên hệ có điều kiện); hướng nghiên cứu tiếp theo\\
\bottomrule
\end{tabularx}}
\nguon{Nguồn: Thí sinh đề xuất}
\end{table}

Thứ tự chương bám theo thứ tự câu hỏi, và mỗi chương thực nghiệm (4, 5, 6) đủ độc lập để thành một bài báo: chương~5 thành bài báo 1, chương~6 thành bài báo 2, chương~4 thành bài báo 3 (mục~7.3). Chương~6 được đặt cuối dù trả lời câu hỏi trọng tâm, vì nó dùng kết quả của chương~5 làm nền trước sự kiện và cần 24 tháng dữ liệu sau sự kiện, nên là chương hoàn thành muộn nhất.
